%\documentclass{beamer}
%\usepacakge[UTF8]{ctex}
%----------------------------------------------------------------------------------------
%    PACKAGES AND THEMES
%----------------------------------------------------------------------------------------

\documentclass[aspectratio=169,xcolor=dvipsnames]{beamer}
\usepackage[UTF8]{ctex}
\usetheme{SimplePlus}
\setbeamertemplate{footline}{}

\usepackage{hyperref}
\usepackage{graphicx} % Allows including images
\usepackage{booktabs} % Allows the use of \toprule, \midrule and \bottomrule in tables

\usepackage{mathpartir}

\usepackage{tikz-cd}
\usetikzlibrary{arrows.meta}
%\tikzcdset{arrow style=tikz, diagrams={>=Stealth}}

\usepackage{fontspec}

\newcommand\calV{\mathcal{V}}
\newcommand\calU{\mathcal{U}}
\newcommand\bfC{\mathbf{C}}
\newcommand\calE{\mathcal{E}}
\newcommand\bbD{\mathbb{D}}
\newcommand\bbI{\mathbb{I}}
\newcommand\bbF{\mathbb{F}}
\newcommand\JEq{\mathtt{JEq}}
\newcommand\FTy{\mathtt{FTy}}
\newcommand\comp{\mathtt{comp}}
\newcommand\fst{\mathtt{fst}}
\newcommand\snd{\mathtt{snd}}
\newcommand\Ty{\mathtt{Ty}}
\newcommand\Ter{\mathtt{Ter}}
\newcommand\Path{\mathtt{Path}}
\newcommand\dM{\mathsf{dM}}
\newcommand\ttp{\mathtt{p}}
\newcommand\ttq{\mathtt{q}}
\newcommand\bfone{\mathbf{1}}
\newcommand\Set{\mathbf{Set}}
\newcommand\DM{\mathbf{DM}}
\newcommand\Id{\mathtt{Id}}
\newcommand\refl{\mathtt{refl}}
\newcommand\C{\mathbf{C}}
\newcommand\psC{\widehat{\mathbf{C}}}
\newcommand\extop{{\scalebox{0.4}[1]{\_}}.{\scalebox{0.8}[1]{\_}}}
\newcommand{\shortminus}{\raisebox{0.3ex}{\scalebox{0.8}[1.0]{-}}}
\newcommand{\mathshortminus}{\mathbin{\text{\normalfont\shortminus}}}
\newcommand{\ul}{{\scalebox{0.4}[0.4]{\_}}}
\newcommand{\ax}[1]{\mathbf{ax_{#1}}}
\newcommand\uc{\mathtt{uc}}

\DeclareMathOperator\FreeDeMorgan{FreeDeMorgan}
\DeclareMathOperator\obj{obj}
\DeclareMathOperator\y{\mathbf{y}}
%\newcommand\cof{\mathtt{cof}}
\newcommand\Cof{\mathtt{Cof}}
\DeclareMathOperator\cof{\mathtt{cof}}

\usepackage{unicode-math}
%\setmathfont{Latin Modern Math}
%\setmathfont{XITS Math}[range=cal]
%\setmathfont{Euler Math}[range=cal]

\DeclareMathAlphabet{\mathcal}{OMS}{cmsy}{m}{n}

\newcommand\cubicalsets{\widehat{□}}

\usepackage{mathpartir}

\newfontfamily{\codefont}{Unifont}
% 设置 verbatim 默认字体
\makeatletter
\renewcommand{\verbatim@font}{\codefont\small}  % 替换为实际的Unifont命令
\makeatother

%----------------------------------------------------------------------------------------
%    TITLE PAGE
%----------------------------------------------------------------------------------------

\title{使用 Agda 形式化}
%\subtitle{Subtitle}

%\author{Pin-Yen Huang}

%\institute
%{
%    Department of Computer Science and Information Engineering \\
%    National Taiwan University % Your institution for the title page
%}
%\date{\today} % Date, can be changed to a custom date
\date{}

%----------------------------------------------------------------------------------------
%    PRESENTATION SLIDES
%----------------------------------------------------------------------------------------

\begin{document}

\begin{frame}
    % Print the title page as the first slide
    \titlepage
\end{frame}

\begin{frame}{背景}
  在拓扑斯的内部类型理论中工作, 而不是从外部处理拓扑斯的对象和态射, 这种方法的一个优点是它更适合机器辅助形式化. 要理解为什么, 可以考虑形式化一个从外部构建的模型会涉及什么. 首先, 我们需要定义所有的数学结构(范畴、函子等)并形式化模型构建中使用的所有相关属性. 接下来, 我们必须形式化模型本身, 使用前一步中定义的结构. 通过在内部语言中工作, 我们基本上可以去掉整个第一步. 这是因为, 我们不是在证明助手中形式化拓扑斯的概念及其内部语言, 而是简单地将证明助手当作拓扑斯的内部语言来使用.
\end{frame}

\begin{frame}[fragile,shrink]{拓扑斯内部语言和 Agda 的区别}
  \begin{itemize}
  \item Agda 的相等类型是内涵性的, 而内部类型理论的相等类型是外延性的. 这意味着我们在形式化过程中需要做更多的工作. 例如, 我们可能需要在命题相等的类型之间进行显式强制转换, 而这在内部语言中是隐式的.
  \item Agda 默认情况下不包含一个与子对象分类子 $Ω$ 对应的非直谓性命题全集.
  \begin{verbatim}

   {-# OPTIONS --type-in-type #-}
   record Ω : Set where
     constructor prop
     field
       prf : Set
       equ : (u v : prf) → u ≡ v
  \end{verbatim}
  我们仅将 \verb|type-in-type| 用于 $Ω$ 的定义, 并且该定义被分离到一个单独的模块中, 而 \verb|type-in-type| 选项在该模块之外不被启用, 我们相信这种受限的使用是相容的.
\end{itemize}
\end{frame}

\begin{frame}[fragile,shrink]{余纤维化}
  考虑一类单态射 $m : A\rightarrowtail B$, 其分类态射
  \[
  \lambda(y:B)\to\exists(x:A). m\,x = y : B \to \Omega
  \]
  可以通过 $\Cof\rightarrowtail\Omega$ 进行分解. 我们称此类单态射为余纤维化.
  \[
  \begin{tikzcd}[row sep = tiny]
  A \arrow[r,tail, "m"] \arrow[dd] & B \arrow[dd] \arrow[dr] & \\
  & & \Cof \arrow[dl,tail] \\
  1 \arrow[r,"\top"] & \Omega &
  \end{tikzcd}
  \]
\end{frame}

\begin{frame}[shrink]{$\ax{7}:[\forall(\varphi\:\psi:\Omega).\cof\varphi\Rightarrow(\varphi\Rightarrow\cof\psi)\Rightarrow\cof(\varphi\land\psi)]$}
  公理 $\ax{7}$ 等价于要求余纤维化在复合运算下封闭.

  \bigskip
  首先假设$\ax{7}$成立, 并且$f:A\rightarrowtail B$和$g:B\rightarrowtail C$ 是余纤维化. 因此
  \begin{itemize}
    \item $\forall(b:B).\cof(\exists(a:A).f\,a=b)$
    \item $\forall(c:C).\cof(\exists(b:B).g\,b=c)$
  \end{itemize}
  都成立, 并且我们希望证明$\forall(c:C).\cof(\exists(a:A).g(f\,a)=c)$. 注意到, 对于 $b:B$ 和 $c:C$
  \begin{align*}
    g\,b=c \quad&\stackrel{g\text{是单态射}}{\implies}\quad (\exists(a:A).g(f\,a)=c) = (\exists(a:A).f\,a=b) \\
    &\stackrel{\phantom{g\text{是单态射}}}{\implies}\quad \cof(\exists(a:A).g(f\,a)=c) = \cof(\exists(a:A).f\,a=b) = \top
  \end{align*}
  因此对于$\varphi\triangleq\exists(b:B).g\,b=c$和$\psi\triangleq\exists(a:A).g(f\,a)=c$, 我们有$\cof\varphi$且$\varphi\Rightarrow\psi$. 因此根据 $\ax{7}$, 我们得到$\cof(\varphi\land\psi)$, 它等于 $\cof\psi$ 因为 $\psi\Rightarrow\varphi$. 所以我们确实有 $\forall(c:C).\cof(\exists(a:A).g(f\,a)=c)$.
\end{frame}
\begin{frame}[shrink]{$\ax{7}:[\forall(\varphi\:\psi:\Omega).\cof\varphi\Rightarrow(\varphi\Rightarrow\cof\psi)\Rightarrow\cof(\varphi\land\psi)]$}
1
\end{frame}


\begin{frame}{唯一选择公理}
  \[
  \uc : (\varphi : A \to \Omega) \to [\exists!(a:A).\varphi\,a]\to\{a:A\mid\varphi\,a\}
  \]
  其中$\exists!(a:A).\varphi\,a\triangleq\exists(a:A).(\varphi\,a\land\forall(a':A).\varphi\,a'\Rightarrow a=a')$.
\end{frame}

\begin{frame}{同构}
  给定对象$A,B:\calU$, $A$和$B$之间的一个同构是一个函数$f:A\to B$, 它有一个双边逆. 令$A\cong B$为$A$和$B$之间的同构类型, 定义为
  \[
  A\cong B\triangleq \{f:A\to B\mid(\exists g:B\to A)(g\circ f=id)\land(f\circ g=id)\}
  \]
  如果存在一个同构$f:A\cong B$, 我们就说$A$和$B$是同构的. 如果两个类型族$A,B:\Gamma\to\calU$的每一个纤维都是同构的, 我们就说它们是同构的, 并且我们不太正规地写成
  \[
  A\cong B \triangleq (x:\Gamma)\to A\,x\cong B\,x
  \]
  同构的逆是相对于内部类型理论的外延等价而言的, 这与我们稍后将介绍的等价概念形成对比, 后者只具有相对于路径等价的逆. 此外, 这些逆是唯一的, 因此, 使用唯一选择, 我们总是可以构造一个同构的逆. 因此, 给定 $f : A ≅ B$, 我们将 $f⁻¹ : B ≅ A$ 写成 $f$ 的这个逆. 
\end{frame}

\begin{frame}[shrink]{严格性公理}
  这条公理关注解决内部语言方法的弱点.

  \bigskip

  任何余纤维-部分类型 $A$, 如果在 $A$ 定义的任何地方都与全类型 $B$ 同构, 则可以扩展为与 $B$ 同构的全类型 $B'$:
  \begin{align*}
    \ax{9} :\ &\{\varphi:\Cof\}(A:[\varphi]\to\calU)(B:\calU)(s:(u:[\varphi])\to(A\,u\cong B)) \to \\
    &(B':\calU)\times\{s':B'\cong B\mid\forall(u:[\varphi]).\ A\,u=B'\land s\,u=s'\}
  \end{align*}
\end{frame}

\end{document}